
\section{模型假设}

为简化问题并聚焦核心分析，本文做出以下假设：

\begin{itemize}[itemindent=2em]
\item \textbf{假设1：}导电颗粒为理想几何体，金属A为严格的直圆柱体（底面半径30nm，长度5000nm），金属B为严格的球体（半径200nm），不考虑表面粗糙度和制造公差。
\item \textbf{假设2：}导电判定仅取决于颗粒表面间最短距离是否不超过隧穿距离1.8nm，不考虑接触电阻、量子效应随距离的连续变化等更复杂的导电机理。
\item \textbf{假设3：}颗粒在微构体内随机均匀分布，圆柱体轴向在各方向等概率分布（球面上均匀采样），不考虑重力沉降、团聚效应或制造过程中的定向排列。
\item \textbf{假设4：}周期性边界条件精确适用，即颗粒超出微构体边界的部分严格按边长平移规则从对侧重新进入，不考虑边界附近的介质不均匀性。
\end{itemize}
